{\begingroup
\newcommand{\rep}{\operatorname{rep}}
\newcommand{\icoh}{\operatorname{IndCoh}}
\newcommand{\qcoh}{\operatorname{QCoh}}
\newcommand{\coh}{\operatorname{Coh}}
\newcommand{\dmod}{\operatorname{D-mod}}
\newcommand{\g}{\mathfrak{g}}
\newcommand{\ug}{U(\mathfrak{g})}
\newcommand{\gmod}{\mathfrak{g}\text{-mod}}
\newcommand{\dr}{\text{dR}}
\newcommand{\spec}{\operatorname{Spec}}

\chapter{Motivation}

\section{Ind-coherent sheaves}
An important method to study modern representation theory is algebro-geometric. Specifically,
For linear algebraic group $G$, \textbf{\color{blue} the category of algebraic representation $\rep(G)$ can be interpreted as 
the category of quasi-coherent sheaves on the classifying stack $BG$.} In this view, many useful functors or operations on $G$ are naturally
correspond to the sheaf level. Explicit formulas do not work well when modern representation theory are inevitably involved in higher category.

The most significant mathematic object in this book is Ind-coherent sheaves. Before we define it, we should know the reason to drop 
the Quasi-coherent sheaves that Lurie used in his higher algebra. 

To see the reason, let $\g$ be the Lie algebra of $G$ with associated formal group $\exp(\g)$.   
Interpret $\gmod$ as modules over $\ug$ and $\rep(G)$ as modules of compactly supported distributions $\mathcal{E}'(G)$.
In this sense, The restriction functor 
\begin{align*}
    \rep(G) \to \gmod
\end{align*}
can be interpreted as the pullback functor along the map
\begin{align*}
    B(\exp(\g)) \to BG
\end{align*}
Since one can prove that the category $\gmod$ is canonically equivalent to $\qcoh(B(\exp(\g)))$. 
Let us now be given a homomorphism $\alpha:\g'\to\g$. The functor of restriction $\gmod\to\mathfrak{g}'\text{-mod}$ is the pullback functor of 
\begin{align*}
    f_\alpha:B(\exp(\g')) \to B(\exp(\g)). \tag{$\bigstar$}
\end{align*}
However the left adjoint to restriction functor is induction functor which cannot correspond to the direct image along $(\bigstar)$.
More explicitly speaking, in the sense of representation theory, the direct image functor as the right adjoint of pullback functor of $(\bigstar)$ is coinduction functor.
In the case of infinite dimension, \textbf{\color{blue}the Frobenius Reciprocity $\text{Ind}\cong\text{Coind}$ fails to hold and thus the geometric realization of induction functor fades.}


The birth of ind-coherent sheaves is aim to deal with the situation that Frobenius Reciprocity fails. For the case of Frobenius Reciprocity holds, for example classical Lie algebra\footnote{Here, "classical" means finite dimensional}, the ind-coherent sheaves are equivalent to quasi-coherent sheaves.

In the sense of ind-coherent sheaves, the restriction functor along the homomorphism $\alpha$ is given by the $!$-pullback functor
\begin{align*}
    (f_\alpha)^!:\icoh(B(\exp(\g)))\to \icoh(B(\exp(\g')))
\end{align*} 
the functor of induction $\g'\text{-mod}\to \gmod$ will be given by the functor of IndCoh direct image 
\begin{align*}
    (f_\alpha)_*^{\text{IndCoh}}:\icoh(B(\exp(\g')))\to \icoh(B(\exp(\g))).
\end{align*}
An analytical analogy is \textbf{\color{blue} $\qcoh$ behaves more like functions on a space, while $\icoh$ behaves more like measures.}

For every scheme of finite type $X$, its de Rham prestack $X_{\text{dR}}$ is a inf-scheme, one can prove that 
\begin{align*}
    \text{IndCoh}(X_{dR}) \cong \text{D-mod}(X)
\end{align*}
Then for an arbitrary inf-scheme $\mathcal{X}$, the ind-coherent sheaf can be given by 
\begin{align*}
    \icoh(\mathcal{X}) = \lim_{Z\to\mathcal{X}}\icoh(Z).
\end{align*}
Hence, one can prove that for any proper map $f:\mathcal{X}'\to\mathcal{X}$, the pair of functor $(f_*^{\icoh},f^!)$ is adjoint.
There are several examples to illustrate its advantages.
\begin{example}
    Take $f:X\to X_{\text{dR}}$ be canonical projection, then the adjoint pair identifies with the pair
    \begin{align*}
        \textbf{ind}_{\text{D-mod}}:\icoh(X)\rightleftarrows \dmod(X): \textbf{ind}_{\dmod}.
    \end{align*}
\end{example}
\begin{example}
    Take $f:Y_{\text{dR}}\to X_{\text{dR}}$, then $f_*^{\icoh}$ and $f^!$ are correspond to the functors
    \begin{align*}
        g_{*,\dmod}:\dmod(Y) \to \dmod(X) \text{\quad and\quad}
        g_{\dmod}^!:\dmod(X) \to \dmod(Y)
    \end{align*}
    of D-module push-forward and pullback respectively.
\end{example}

\begin{example}
    Consider Lie algebra $\g$ acts on a scheme $X$. In this case, we have a diagram 
    \begin{align*}
        B(\exp(\g))\xleftarrow{f_1} B_X(\exp(\g)) \xrightarrow{f_2} X_\dr.
    \end{align*}
    Then the composite functor 
    \begin{align*}
        (f_2)^{\icoh}_*\circ (f_1)^!:\icoh(B(\exp(\g)))\to \icoh(X_\dr)
    \end{align*}
    identifies with the localization functor $\gmod\to \dmod(X)$. 
\end{example}

\begin{remark}
    $BG$ manifests as a geometric representable functor 
    \begin{align*}
        BG: \mathrm{dSch}^{\mathrm{op}} \to \infty\text{-}\mathrm{Grpd} \quad X\mapsto \text{Bun}_G(X).
    \end{align*}
    From an operadic perspective, 
    \begin{align*}
        \text{Bun}_G(X) \simeq \text{Alg}_{X}(BG) \simeq  \mathrm{Map}_{\mathrm{PreStk}}(X, BG)
    \end{align*}
    where we interpret $X$ as an unary $\infty$-operad and $\text{Bun}_G(X)$ as the category of $BG$-valued algebras over $X_{\mathrm{dR}}$.
    When we consider the de Rham prestack $X_\dr$, we obtain that
    \begin{align*}
         \text{LocSys}_G(X) \simeq \text{Alg}_{X_\dr}(BG) \simeq  \mathrm{Map}_{\mathrm{PreStk}}(X_{\text{dR}}, BG).
    \end{align*}
    $B_X(G)$  is the $\infty$-categorical relative tensor product $X\otimes_{G} *$.
\end{remark}
\section{Construction of ind-coherent sheaves}
There is a heuristic formula we have presented 
\begin{align*}
    \icoh(\mathcal{X}) = \lim_{Z\to\mathcal{X}}\icoh(Z) \tag{$\bigstar$}
\end{align*}
In what sense is this limit taken? In triangulated categories, colimits and limits are not well-behaved. 

We do it in the $\infty$-category of presentable $k$-linear stable $\infty$-categories, or equivalently, in the $\infty$-category of differential graded categories:
categories enriched over the category of chain complexes of $k$-modules, such that the differential satisfies the Leibniz rule with respect to composition.
We denote it by DGCat.

For a scheme $Z$, the category $\icoh(Z)$ is the renormalization of $\qcoh(Z)$. 
One can prove that ind-coherent sheaves have the following properties:
\begin{itemize}
    \item[a.] DG category $\icoh(Z)$ contains all infinite coproducts
    \item[b.] For a map $Z'\xrightarrow{f}Z$, the functor $\icoh(Z')\xrightarrow{f^!}\icoh(Z)$ preserves all infinite coproducts.
\end{itemize}
We can see that the shriek pullback functor $f^{!,\qcoh}$ fails to satisfy property (b) even in the case that $f$ is proper.
This is because the direct image functor $f_*$ as the left adjoint of $f^{!,\qcoh}$ does not preserve compactness, and thus its right adjoint does not preserve infinite coproducts according to Brown representability. Indeed, $f_*$ does not necessarily
send perfect complexes on $Z'$ to perfect complexes on $Z$ unless $f$ is of finite Tor-dimension.
An easy example is $f:\spec(k)\to \spec(k[t]/t^2)$. The perfect complex $k$ does not has finit resolution as a $k[t]/t^2$-module.

Fortunately, when $f$ is proper, the direct image functor $f_*$ sends the category of bounded coherent complexes $\coh^b(Z')$ to $\coh^b(Z)$ and we define ind-coherent sheaves as the ind-completion of $\coh^b(-)$
\begin{align*}
    \mathrm{IndCoh}(-) := \mathrm{Ind}(\mathrm{Coh}^b(-))
\end{align*}
The compact objects in the category $\icoh(Z)$ are exactly all bounded coherent complexes, and therefore $f_*$ is well-behaved.



\endgroup}